% =====================================================================
%  CHƯƠNG 10. TỔNG KẾT VÀ ĐÁNH GIÁ HỌC THUẬT
% =====================================================================
\chapter{Tổng kết và đánh giá học thuật}

\section{Đánh giá tổng quan về năng lực và giới hạn hiện tại của AI Agent}
Tổng hợp từ Google Cloud \cite{google_agents}, IBM Think \cite{ibm_agents}, McKinsey \cite{mckinsey_agents} và
GeeksforGeeks \cite{gfg_rl}, có thể khẳng định AI Agent là bước chuyển quan trọng của trí tuệ nhân tạo: từ vai trò trợ lý
thụ động trả lời câu hỏi sang thực thể tự hành giải quyết công việc thực tế. Ba trụ cột suy luận, bộ nhớ và công cụ cho
phép tác tử lập kế hoạch, ghi nhớ và tác động vào hệ thống; các mô hình suy luận như ReAct, ReWOO, Reflection quyết định
sự cân bằng giữa độ linh hoạt và chi phí; mạng lưới tác tử (Agentic Mesh) mở rộng mô hình này ra toàn doanh nghiệp.

Tuy vậy, các thách thức về độ tin cậy, chi phí token, an toàn quyền hạn (privilege escalation) và nguy cơ vòng lặp vô
hạn đòi hỏi các rào chắn nghiêm ngặt: nhật ký kiểm toán, định danh tác tử, ngưỡng giới hạn (guardrails) và cơ chế con
người can thiệp (human-in-the-loop). Mức tự trị của tác tử càng cao thì các cơ chế này càng cần thiết.

\section{Bài học kinh nghiệm từ bài thực hành lập trình Q-Learning}
Qua quá trình lập trình và chạy tác tử Q-Learning trong mê cung
$\SL{ch8_qlearning/params/n_rows} \times \SL{ch8_qlearning/params/n_cols}$:
\begin{itemize}
  \item \textbf{Sức mạnh của cập nhật Bellman.} Tác tử không biết bản đồ ban đầu, nhưng sau
        \SL{ch8_qlearning/params/episodes} episode tương tác (\SL{ch8_qlearning/result/seconds} giây trên CPU) đã tự tìm
        ra lộ trình \SL{ch8_qlearning/result/path_len} bước, đúng bằng đường ngắn nhất do BFS kiểm chứng. Các giá trị
        Q-Table hội tụ về đúng tổng phần thưởng chiết khấu mà phương trình Bellman dự đoán.
  \item \textbf{Tầm quan trọng của cân bằng thăm dò -- khai thác.} Giai đoạn đầu với $\varepsilon$ lớn có phần thưởng
        rất thấp nhưng cần thiết, vì nhờ đó tác tử tình cờ tìm thấy đích và lan giá trị ngược về. Nếu $\varepsilon$ giảm
        quá nhanh, tác tử dễ mắc kẹt ở một lộ trình chưa tối ưu; nếu $\varepsilon$ luôn lớn, chính sách không thể ổn
        định.
  \item \textbf{Thiết kế phần thưởng định hình hành vi.} Hình phạt $-1$ mỗi bước khiến tác tử ưu tiên đường ngắn; hình
        phạt $-10$ khi đâm tường dạy nó tránh tường lửa. Đây chính là vai trò của ``tiêu chuẩn hiệu năng'' trong kiến
        trúc Learning Agent.
  \item \textbf{Guardrail cũng cần cho tác tử nhỏ.} Giới hạn \SL{ch8_qlearning/params/max_steps} bước mỗi episode ngăn
        tác tử lặp vô hạn ở giai đoạn đi ngẫu nhiên, đúng như nguyên tắc \texttt{max\_iterations} mà IBM khuyến nghị.
\end{itemize}

\section{Hướng phát triển trong các bài tập kế tiếp}
Mở rộng kiến trúc tác tử kết hợp đồ thị tri thức (Knowledge Graph RAG) và tích hợp công cụ Function Calling cho các bài
toán tài chính và y tế phức hợp trong các bài thực hành tiếp theo của học phần.
